\paragraph{Writing a character's meaning in context.}
A new 1.18-million-parameter writer learned to draw English glosses for a marked character in classical Chinese. It received the character and up to eight neighbouring cells on either side. We trained the entire model from random initialization on 247,456 annotated occurrences from six works. No English letters appeared in the input passages. On the held portions of the \textit{Analects} and \textit{Chu ci}, 5,923 of 6,987 and 5,808 of 7,173 emitted images matched the treebank gloss exactly, including the learned ending. Many headwords had a dominant reading. A count predictor using local contexts recovered 6,306 and 6,528 glosses.

The more discriminating test required two different glosses for the same printed character in two passages (Fig.~\ref{fig:visual-word-senses}). The aligned writer solved both passages for 213/566 and 82/403 pairs. Permuting the glosses among training occurrences of the same headword gave 4/566 and 5/403. Paired 95\% intervals for the differences, resampled by headword, were 31.6--42.4 and 12.6--25.3 percentage points. The count predictor solved 211 and 104 pairs. For the character meaning speech, the writer completed 15/38 pairs in the \textit{Analects} and 4/8 in \textit{Chu ci}. Its autonomous replies included the images of \textit{speak} and \textit{speech}. None of 69 instances requiring a gloss string unseen during fitting was exact. When both candidates were provided for a separate 128 passages, the visual writer selected 80 reference glosses and text Qwen3-8B selected 79. They completed 18 and 20 of 64 pairs respectively.

Matching the treebank can overstate understanding. With the same three central cells in both passages, the aligned writer solved 2/27 contrasts in the \textit{Analects} and 0/11 in \textit{Chu ci}. With five fixed cells, it solved one of six and zero of three. Its sole five-cell success involved the same phrase about rites, music and military expeditions in two parallel clauses (\textit{Analects} 16.2). The treebank glosses the character for music as \textit{music} in one and \textit{enjoy} in the other. Legge's translation reads the two clauses in parallel. We retain the reference score and give this pair no credit as evidence for a change of meaning.

The aligned language fit used 38.8 s, 5.51 kJ and 326 MB peak allocation on one RTX 4090 D. The shared glyph drawing network took a further 0.41 s. An unseen serif input font reduced exact replies to 1,458 and 1,157 on the two works. Mapping source cells to the nearest training ink before reading raised these counts to 5,429 and 5,410. In a second unseen printed face resembling brush script, the same procedure raised exact replies from 1 to 1,209 of 6,987 and from 0 to 1,091 of 7,147 font-covered passages. The remaining gap measures sensitivity to how a word is printed, even with the source words and requested reading fixed.
